\documentclass[../../../include/open-logic-section]{subfiles}

\begin{document}

\olfileid{sth}{z}{infinity-again}	
\olsection{نامتناهيت}

تر اوسه دومره بديهي اصلونه لرو چې د بې‌نهايته ډېرو سټونو شتون
تضمين کړي (که کوم سټونه وي). ځکه فرض کړئ چې کوم سټ شتون لري؛
نو $\emptyset$ شتون لري (د \olref[sth][z][sep]{prop:emptyexists} له مخې).
اوس د هر سټ~$x$ دپاره سټ $x \cup \{x\}$ د
\olref[sth][z][pairs]{prop:pairsconsequences} له مخې شتون لري.
نو د دې اصل په يو څو ځله کارولو سره به لاندې سټونه ترلاسه کړو:
\begin{enumerate}
	\item[0.] $\emptyset$ %& $0$ & stage $0$\\
	\item[1.] $ \{\emptyset\}$ %& $1$ & stage $1$\\
	\item[2.] $\{\emptyset, \{\emptyset\}\}$ %& $2$ & stage $2$\\
	\item[3.] $\{\emptyset, \{\emptyset\}, \{\emptyset, \{\emptyset\}\}\}$ %&$3$ & stage $3$\\
	\item[4.] $\{\emptyset, \{\emptyset\}, \{\emptyset, \{\emptyset\}\}, \{\emptyset, \{\emptyset\}, \{\emptyset, \{\emptyset\}\}\}\}$% & $4$ & stage $4$ 
\end{enumerate}
او کتلے شو چې دا سټونه ټول له يو بل څخه بېل دي.

د څو وجوهاتو له مخې مو شمېرنه له $0$ څخه پېل کړې ده. خو يو
سبب ئې دا دے. دا کتل ډېر ګران نۀ دي چې هغه سټ چې «$n$» نوم مو
ورکړے دے، په کره توګه $n$ غړي لري او (په شهودي توګه) په
$n$م پړاو کښې جوړېږي.

خو دا \emph{بې‌نهايته ډېر} سټونه راکوي، مګر د يوه \emph{نامتناهي سټ}
شتون نۀ تضمينوي؛ يعنې د داسې سټ شتون چې بې‌نهايته ډېر غړي ولري.
او دا خبره ډېره مهمه ده: تر څو د ډېډېکېنډ له مخې يو نامتناهي سټ
ونۀ مومو، د ډېډېکېنډ الجبر نۀ شو جوړولے. خو د ډېډېکېنډ الجبر
غواړو، څو ئې د طبيعي عددونو د سټ په توګه واخلو. (له
\olref[sfr][infinite][dedekindsproof]{sec} سره پرتله کړئ.)

هغه بديهي اصلونه چې تر اوسه مو ټاکلي دي، د هېڅ نامتناهي سټ
شتون \emph{نۀ} تضمينوي. نو يو نوے بديهي اصل ټاکل پکار دي:

\begin{axiom}[نامتناهيت]
يو سټ $I$ شته چې $\emptyset \in I$ وي او $x \cup \{x\} \in I$ وي، هر کله چې $x \in I$ وي.
\begin{align*}
	\exists I( & (\exists o \in I)\forall x\ x \notin o \land {}\\
	& (\forall x \in I)(\exists s \in I)\forall z(z \in s \liff (z \in x \lor z = x)))
\end{align*}
\end{axiom}

په اسانۍ ليدلے شو چې هغه سټ $I$ چې د نامتناهيت بديهي اصل ئې
راکړي، د ډېډېکېنډ له مخې نامتناهي دے. د هغه ځانګړے غړے
$\emptyset$ دے او په $I$ باندې يو پر يو تابع د $s(x) = x\cup
\{x\}$ په وسيله ورکول کېږي. اوس
\olref[sfr][infinite][dedekind]{thm:DedekindInfiniteAlgebra} وښودله
چې د ډېډېکېنډ له مخې له يوه نامتناهي سټ څخه د ډېډېکېنډ الجبر
څنګه راايستلے شو؛ او دا به د طبيعي عددونو د خپل سټ په توګه واخلو.
په لا دقيقه توګه: \emph{د سرچينې د حدودو يادښت (OLSTH-003): ليکل
شوے نامتناهيت اصل په هر راکړل شوي شاهد کښې د همدې تالي تابعې
يو پر يو توب جلا نۀ ثابتوي. دلته د تجمعي پړاوونو غير رسمي تصور
کارېدلے دے؛ د طبيعي عددونو لږ تر لږه بنده برخه بايد د هر شاهد له
ټولو غړو سره يو نۀ وګڼل شي. د بنياد يا په پړاوونو کښې د جوړېدنې
اضافي فرض بې له څرګندونې د همدې اصل برخه نۀ ده.}
\begin{defn}\ollabel{defnomega}
$I$ هغه هر سټ وګرځوئ چې د نامتناهيت بديهي اصل ئې راکړي.
$s$ تابع $s(x) = x \cup \{x\}$ وګرځوئ.
$\omega =
\closureofunder{s}{\emptyset}$ وګرځوئ. د $\omega$ غړو ته
\emph{طبيعي عددونه} وايو، او وايو چې $n$ د $n$ ځله د $s$ د
کارولو پايله ده، چې پر $\emptyset$ باندې ترسره شوي دي.
\end{defn}

اوس بېرته کتلے شئ چې هغه سټ چې يو څو پراګرافه مخکې «$n$»
نوم ورکړل شوے و، د عدد $n$ \emph{په توګه} به اخيستل کېږي.

د دې ټاکنې د اهميت په اړه به په \olref[sth][z][nat]{sec} کښې
بحث وکړو. د اوس دپاره دا موږ ته د يوې شهودي پايلې ثبوت ممکنوي:

\begin{prop}\ollabel{naturalnumbersarentinfinite}
هېڅ طبيعي عدد د ډېډېکېنډ له مخې نامتناهي نۀ دے.
\end{prop}

\begin{proof}
ثبوت د استقرا په وسيله دے، يعنې د
\olref[sfr][infinite][induction]{thm:dedinfiniteinduction} له مخې.
څرګنده ده چې $0
= \emptyset$ د ډېډېکېنډ له مخې نامتناهي نۀ دے. د استقرا د ګام
دپاره به نقيض عکس ثابت کړو: که (په نامعقوله توګه) $s(n)$ د
ډېډېکېنډ له مخې نامتناهي وي، نو $n$ هم د ډېډېکېنډ له مخې نامتناهي دے.

نو فرض کړئ چې $s(n)$ د ډېډېکېنډ له مخې نامتناهي دے، يعنې
يوه !!{injection} $f$ شته چې $\ran{f}\subsetneq \dom{f} = s(n) = n \cup
\{n\}$ وي. دوه حالتونه په پام کښې نيول پکار دي.

\emph{لومړے حالت: $n \notin \ran{f}$.} نو $\ran{f} \subseteq n$ او
$f(n) \in n$ دي. $g = \funrestrictionto{f}{n}$ وګرځوئ؛ اوس $\ran{g} =
\ran{f} \setminus \{f(n)\} \subsetneq n = \dom{g}$ لرو. نو $n$
د ډېډېکېنډ له مخې نامتناهي دے.

\emph{دويم حالت: $n \in \ran{f}$.} يو $m \in \dom{f} \setminus \ran{f}$
وټاکئ، او تابع $h$ د تعريف ساحې $s(n) = n \cup \{n\}$ سره تعريف کړئ:
\[
h(x) =
	\begin{cases}
		f(x) & \text{که }f(x) \neq n\\
		m & \text{که }f(x)=n
	\end{cases}
\]
نو $h$ او $f$ په هر ځاے کښې يو شان دي، پرته له دې چې $h(f^{-1}(n)) = m
\neq n = f(f^{-1}(n))$ دي. ځکه چې $f$ !!a{injection} ده، نو $n \notin
\ran{h}$؛ او $\ran{h} \subsetneq \dom{h} = s(n)$ دي. اوس $n$
د لومړي حالت د استدلال په کارولو سره د ډېډېکېنډ له مخې نامتناهي دے.
\end{proof}

خو دا پوښتنه لا پاتې ده چې د نامتناهيت بديهي اصل څنګه
\emph{توجيه} کولے شو. لنډ ځواب دا دے چې هغې کيسې ته چې بيان کړې
مو ده، يو بل اصل زياتول پکار دي. هغه اصل دا دے:
\begin{enumerate}
	\item[] \stagesinf. يو نامتناهي پړاو شته. يعنې داسې يو پړاو شته
چې (a) لومړنے پړاو نۀ دے، او (b) ترې مخکې ځينې پړاوونه شته،
خو (c) سمدستي مخکښ پړاو نۀ لري.
\end{enumerate}
له دې اصل څخه د نامتناهيت بديهي اصل په نېغه نتيجه کېږي.
پوهېږو چې طبيعي عدد $n$ په پړاو $n$ کښې جوړېږي. نو سټ
$\omega$ په لومړي نامتناهي پړاو کښې جوړېږي. او همدا $\omega$
د نامتناهيت بديهي اصل شاهد دے.

خو دا موږ بېرته هغې پوښتنې ته ستنوي چې \stagesinf څنګه توجيه کولے
شو. د \stagessucc په شان، دا هم د تجمعي-تکراري مراتبو په اړه زموږ
د کيسې ښکاره برخه نۀ وه. خو خبره تر دې هم وړاندې ده: د تکراري
مراتبو په لړۍ کښې، چې سټونه پړاو په پړاو جوړېږي، هېڅ شے دا
نۀ لازموي چې په دې بهير کښې يو \emph{نامتناهي} پړاو هم ومنو.
دا تصور پوره سازګار ښکاري چې پړاوونه د طبيعي عددونو په شان مرتب وي.

خو له دې نه يوه څرګنده ستونزه پيدا کېږي. په
\olref[sfr][infinite][dedekindsproof]{sec} کښې مو د ډېډېکېنډ هغه
«ثبوت» وڅېړه چې د (فکرونو) يو د ډېډېکېنډ له مخې نامتناهي سټ شته.
کېدلے شي هغه ثبوت مو چندان قانع کوونکے نۀ وي ګڼلے. خو که
\stagesinf{} د سټ د تکراري تصور (يا «د فکر د قوانينو») له مخې
«راپورې نۀ تپل کېږي»، نو لا هم د ډېډېکېنډ له مخې د يوه نامتناهي
سټ د شتون د ادعا ذاتي توجيه نۀ لرو.

بېشکه دلته ډېر نور څه هم ويل کېدلے شي. خو هيله ده اوس دې ځاے ته
رسېدلي يئ چې دا فکر پېل کړئ چې د يوه بديهي اصل (يا اصل) د
توجيه دپاره \emph{څه پکار} دي. په راتلونکي بحث کښې به
\stagesinf{} يوازې منلے فرض کړو.

\end{document}
